\section{Six-player Canadian Fish}
\label{sec:game}

This section states the game completely, so that nothing later in the paper depends on outside
material. Readers who know the game should still read \S\ref{sec:game-dialect}: published
descriptions differ on several consequential points, and the dialect studied here is one explicit
selection among them.

\subsection{The rules}
\label{sec:game-rules}

Canadian Fish --- also called Literature --- is played by six players in two teams of three, seated
alternately so that no two teammates are adjacent. The deck is 54 cards: the standard 52 plus two
jokers. It is partitioned into nine \textbf{half-suits} of six cards each --- for each of the four
suits, the low half $2$--$7$ and the high half $9$--$A$ --- plus a ninth half-suit made of the four
eights and the two jokers. Every card belongs to exactly one half-suit, and each player is dealt
nine cards. The team that takes at least five of the nine half-suits wins.

Play proceeds by \textbf{asks}. On their turn a player names one opponent and one card. The ask is
legal only if

\begin{enumerate}
\item the named card lies in a half-suit that has not yet been claimed;
\item the asker holds at least one \emph{other} card of that half-suit;
\item the asker does \emph{not} hold the named card; and
\item the target is an opponent who still holds at least one card.
\end{enumerate}

If the target holds the card they must surrender it, and the asker keeps the turn. If they do not,
the turn passes to the target. Condition~(2) is what makes the game a game: a player may only ask
about a half-suit they already have a stake in, so every ask is a public declaration of partial
ownership.

A half-suit is claimed by a \textbf{declaration}. A declaring player names, for each of the six
cards of a half-suit, which member of their \emph{own} team holds it. The claim is all-or-nothing:
if every one of the six assignments is correct the declaring team scores the half-suit, and if any
assignment is wrong the opposing team scores it. A player may not declare a half-suit they hold
entirely --- the interesting case is always a claim about cards in teammates' hands.

\textbf{Everything is public.} Every ask, its answer, every card transfer, every declaration and its
outcome, and every player's card count are observed by all six players. Nothing is hidden except the
initial deal.

\subsection{The dialect studied, and why the choice is explicit}
\label{sec:game-dialect}

Published descriptions of this game document mutually incompatible variants \cite{pagat}, so the
rule set below is a selection rather than a canon. It matches the strategy corpus this line of work
draws on \cite{develin} and the rules of record supplied by the project owner. Every contested
choice is an engine switch, so its effect is measurable rather than assumed, and
\S\ref{sec:results-dialects} measures the headline result under \vsevenDiaRows{} of them.

The consequential choices are:

\begin{itemize}
\item a 54-card deck in nine half-suits, nine cards per player;
\item declaration legal at any moment, in or out of turn, by a seat holding cards or not;
\item a misdeclaration awarding the half-suit to the opposing team;
\item lowest seat index first among simultaneous declarers;
\item an eight-level willingness ladder in the \textbf{forced endgame}, which arises when one whole
      team is cardless and the other must continue declaring;
\item a 400-ask action cap, whose residue is adjudicated neutrally by physical majority.
\end{itemize}

No prearranged signalling conventions are available to any policy. Team communication is limited to
the forced-endgame willingness bit and to the successor a cardless turn-holder names.
Appendix~\ref{app:dialect} lists each choice against the switch that changes it.

\subsection{Two places the engine is deliberately more restrictive than the rules allow}
\label{sec:game-restrictive}

Both are coordination channels a policy could in principle exploit, and both were measured rather
than argued about. They are stated here because a reader entitled to ask ``why not let the team talk
about it?'' deserves the number rather than the principle.

\paragraph{Turn transfer.} When the turn reaches a cardless seat whose teammates both still hold
cards, the rules of record permit the team to negotiate openly over who receives it, sharing nothing
but willingness. The engine instead has the cardless seat choose unilaterally from its own belief.
The channel is empty: a genuine multi-candidate transfer decision arises \numTransferEvents{} times
per game, and replacing the unilateral rule with a \emph{ground-truth} chooser on one team, paired on
common deals, is worth $\numTransferDelta \pm \numTransferCI$ half-suits per game over
\numTransferGames{} paired games, with the two arms diverging at all in \numTransferDiverged\% of
games. A legal mechanism cannot beat the clairvoyant one, so that interval bounds the entire
channel.

\paragraph{Declaration arbitration.} The rules do not say who acts first when two seats offer a
declaration at the same instant, so the order is a modelling choice. The engine scans by lowest seat
index, deliberately declining to rank candidates by stated confidence, because a confidence
comparison moves a statistic of one seat's hand to five others. That refusal costs
\numArbCost{}~pp of win rate pooled over \numArbGames{} games (range \numArbRange{}~pp across five
seeds), and a clairvoyant arbitrator --- one that picks a proposer whose named allocation is actually
correct --- is worth \numArbOracle{}~pp on the same design. The information-safe rule forgoes about a
third of a percentage point, and no rule of any kind, safe or unsafe, informed or omniscient,
recovers appreciably more.
